\documentclass[../main.tex]{subfiles}
\begin{document}

\chapter{Memory Ordering}
\lessoninfo{
  video={Lesson 9, videos 1--14},
  textbook={H\&P \cite{hennessy2017} §3.5--3.6},
  keywords={memory dependence, load, store, commit, load-store queue, LSQ, store-to-load forwarding, out-of-order load-store execution, in-order load-store execution, re-execution, reorder buffer, reservation station},
}

With the reorder buffer of Lesson 8, an out-of-order processor handles every
dependence that passes through registers and keeps exceptions precise.  Loads
and stores, however, also depend on one another through memory, and whether a
load reads the location that an earlier store writes is known only once both
addresses have been computed, which is why Lesson 7 performed loads and stores
in program order.  This lesson shows how they can execute out of order: stores
write memory only at commit, a load-store queue keeps the uncommitted loads and
stores in program order, and each load obtains its value either from memory or
directly from the store that wrote it.

\section{Dependences through memory}
\label{sec:l09-dependences}

The processor of Lesson 8 deals with control dependences and with every
dependence that passes through a register.\source{Lesson 9, videos 1--2; H\&P
§3.5--3.6.}  Branch prediction
(Lesson 4) removes most of the cost of control dependences, and the ROB cancels
the instructions of a mispredicted path; renaming (Lesson 6) removes the false
dependences between writes to the same register; the scheduling of Tomasulo's
algorithm (Lesson 7) makes each instruction wait until the register values it
reads have been computed; and the ROB makes exceptions precise.

Loads and stores remain.  They depend on one another through memory locations
in the same way as instructions depend on one another through registers
(\cref{sec:l07-load-store}):
\begin{lstlisting}[language={[x86masm]Assembler}, numbers=none]
SW  R1, 0(R2)      ; MEM[R2] <- R1
ADD R4, R5, R6     ; R4 <- R5 + R6
LW  R3, 0(R4)      ; R3 <- MEM[R4]
\end{lstlisting}
The load has a RAW dependence on the store exactly when R4 equals R2.  Two
properties distinguish such a dependence from a dependence through a register.
\begin{itemize}
  \item It is not visible in the instructions themselves.  The register names
    say nothing about the memory location, which is known only after the
    address has been computed, often several cycles after the instruction was
    issued.
  \item Nothing renames memory.  The RAT maps register names only, so the
    processor must find, for every load, the store whose value the load has to
    read, or determine that there is none.
\end{itemize}
The rest of the lesson answers two questions: when a store actually writes
memory, and where a load finds its value.

\section{When memory is written}
\label{sec:l09-write}

A store executes like any other instruction once its operands are available,
but it does not write memory at that time.\source{Lesson 9, video 3; H\&P
§3.6.}  Executing the store computes the address and determines the value; the
write itself takes place when the store commits.  A write to memory can no more
be undone than a write to a register.  If a store wrote memory when it
executed, an exception raised afterwards by an earlier instruction would let
the handler see a location that, in program order, has not yet been written,
and a store on a mispredicted path would change memory that the program never
changes.  At commit, in contrast, the store is known to belong to the
execution of the program, and stores commit in program order.

Deferring the write creates a problem for loads.  A load could read memory as
soon as it has computed its address, but a store that precedes it in program
order may have executed and still be waiting to commit, or may not even have
computed its address yet.  If that store writes the location the load reads,
memory still holds the old value, and the load must instead obtain the value of
the store.\margin{In the pipeline of Lesson 3, the MEM stage performs loads and
stores in program order.}  The processor must therefore keep, in program order,
the address and value of every store that has not yet committed, and every load
must consult them.

\begin{keypoint}
  Stores execute out of order but write memory at commit, in program order.
  Until then their values exist only inside the processor, and loads must look
  for them there.
\end{keypoint}

\section{The load-store queue}
\label{sec:l09-lsq}

The uncommitted loads and stores are kept in a queue that works like the
ROB.\source{Lesson 9, videos 4--5; H\&P §3.6.}

\begin{definition}[Load-store queue]
  \label{def:l09-lsq}
  The \term{load-store queue}[ロードストアキュー]\index{LSQ|see{load-store
  queue}} (LSQ) holds, in program order, every load and store that has been
  issued but has not yet committed.  Each entry has four fields:
  \begin{description}
    \item[L/S] whether the instruction is a load or a store;
    \item[Addr] the memory address, once it has been computed;
    \item[Value] for a store, the value to be written; for a load, the value
      it has obtained;
    \item[C] set when the entry is complete: a load has obtained its value, or
      a store has recorded both its address and its value.
  \end{description}
\end{definition}

Like the ROB, the LSQ is a circular buffer with an issue pointer and a commit
pointer (\cref{fig:l09-lsq}).  A load or store receives the entry at the issue
pointer when it is issued, in addition to its ROB entry, and its LSQ entry is
freed when it commits.  Other instructions do not use the LSQ.  A load records
its address and then its value.  A store records its address as soon as its
base register is available and its value as soon as the value register is,
in either order.  Because a store keeps its entry until it commits and writes
memory only then, the LSQ contains exactly the stores whose values are not yet
in memory.\margin{The ROB also holds the loads and stores, but a separate queue
keeps the address searches short.}

\begin{figure}[htbp]
  \centering
  % The load-store queue: loads and stores in program order between the commit
% pointer and the issue pointer.  The load in LSQ6 searches the earlier
% entries, skips loads and non-matching stores, and takes the value of the
% nearest store with the same address.
\begin{tikzpicture}[x=1cm,y=1cm, font=\footnotesize\sffamily, >={Stealth[length=1.8mm]},
  hd/.style={font=\footnotesize\bfseries\sffamily},
  used/.style={draw, fill=ln-accent!12, minimum height=0.45cm, inner sep=0pt},
  free/.style={draw, fill=white, minimum height=0.45cm, inner sep=0pt},
  me/.style={draw, fill=ln-accent!35, minimum height=0.45cm, inner sep=0pt},
  lab/.style={font=\scriptsize\sffamily, text=ln-muted},
  ann/.style={font=\scriptsize\sffamily, anchor=west}]
  % column headings
  \node[hd] at (0.35,0.45) {L/S};
  \node[hd] at (1.25,0.45) {Addr};
  \node[hd] at (2.25,0.45) {Value};
  \node[hd] at (2.95,0.45) {C};
  % rows: entry / type / address / value / style
  \foreach \r/\t/\a/\v/\st in {1/{}/{}/{}/free, 2/S/0x100/17/used,
      3/L/0x108/4/used, 4/S/0x108/25/used, 5/L/{}/{}/used,
      6/L/0x100/{}/me, 7/S/{}/{}/used, 8/{}/{}/{}/free} {
    \pgfmathsetmacro{\y}{-(\r-1)*0.45}
    \node[\st, minimum width=0.7cm] at (0.35,\y) {\t};
    \node[\st, minimum width=1.1cm, font=\scriptsize\ttfamily] at (1.25,\y) {\a};
    \node[\st, minimum width=0.9cm] at (2.25,\y) {\v};
    \node[\st, minimum width=0.5cm] at (2.95,\y) {};
    \node[lab, anchor=east] at (-0.05,\y) {LSQ\r};
  }
  % completion bits: set for LSQ2-LSQ4, and for LSQ6 by the forwarding
  \foreach \y in {-0.45,-0.9,-1.35} {
    \draw[thick] (2.83,\y) -- (2.92,\y-0.09) -- (3.08,\y+0.1);
  }
  \draw[thick, ln-accent] (2.83,-2.25) -- (2.92,-2.34) -- (3.08,-2.15);
  % forwarded value in the searching load
  \node[text=ln-accent, font=\footnotesize\bfseries\sffamily] at (2.25,-2.25) {17};
  % pointers
  \draw[->, thick, ln-accent] (-1.45,-0.45) -- (-0.85,-0.45);
  \node[font=\scriptsize\sffamily, text=ln-accent, anchor=east] at (-1.45,-0.45)
    {\iflangja{コミット}{commit}};
  \draw[->, thick, ln-accent] (-1.45,-3.15) -- (-0.85,-3.15);
  \node[font=\scriptsize\sffamily, text=ln-accent, anchor=east] at (-1.45,-3.15)
    {\iflangja{発行}{issue}};
  % search of the load in LSQ6
  \draw[thick, ln-accent] (3.25,-2.25) -- (3.6,-2.25) -- (3.6,-0.45);
  \draw[->, thick, ln-accent] (3.6,-0.45) -- (3.25,-0.45);
  \foreach \y in {-0.9,-1.35,-1.8} { \draw[ln-accent] (3.52,\y) -- (3.68,\y); }
  \node[ann, text=ln-accent] at (3.8,-2.25)
    {\iflangja{アドレス 0x100 のロード：前方を探索}{load at 0x100: searches earlier entries}};
  \node[ann, text=ln-muted] at (3.8,-1.8) {\iflangja{ロード：無視}{load: skipped}};
  \node[ann, text=ln-muted] at (3.8,-1.35) {\iflangja{ストア，0x108：不一致}{store, 0x108: no match}};
  \node[ann, text=ln-muted] at (3.8,-0.9) {\iflangja{ロード：無視}{load: skipped}};
  \node[ann, text=ln-accent] at (3.8,-0.45)
    {\iflangja{ストア，0x100：一致，値 17 を転送}{store, 0x100: match, value 17 forwarded}};
  \node[ann, text=ln-muted] at (3.8,-2.7)
    {\iflangja{後続のストア：探索しない}{later store: not examined}};
  % program order
  \draw[->, ln-muted] (-2.55,0.2) -- (-2.55,-3.35);
  \node[lab, rotate=90] at (-2.8,-1.57) {\iflangja{プログラム順}{program order}};
\end{tikzpicture}

  \caption{A load-store queue with eight entries; LSQ2 to LSQ7 are in use.  The
    load in LSQ6 has computed the address \texttt{0x100}.  Its search skips
    loads and the store to \texttt{0x108} and takes the value 17 from the
    nearest earlier store to \texttt{0x100}.}
  \label{fig:l09-lsq}
\end{figure}

When a load has computed its address, it records the address in its LSQ entry
and searches the earlier entries, from its own entry back towards the commit
pointer.  Loads are skipped, because reading a location does not change it;
only stores matter.  The search ends at the first store with the same address
or at the commit pointer.
\begin{enumerate}
  \item If it finds a store with the same address, this nearest matching store
    is the last write to the location before the load in program order, so its
    value is the one the load must read.  The load copies the value from the
    store's LSQ entry, waiting for it if the store has not recorded its value
    yet, and does not access memory at all.  Supplying a load with the value of
    an uncommitted store in this way is called
    \term{store-to-load forwarding}[ストアロードフォワーディング].
  \item If it reaches the commit pointer without finding a matching store or a
    store with an unknown address, no uncommitted store writes the location
    before the load, so memory already holds the correct value, and the load
    reads memory.
\end{enumerate}
If, before it ends, the search passes a store whose address has not been
computed yet, that store may or may not write the location the load reads, and
the processor cannot yet tell which.  This case is the difficult one, and the
next section presents the ways of dealing with it.

\section{Out-of-order and in-order load-store execution}
\label{sec:l09-policies}

For a load whose search passes a store with an unknown address, there are
three options.\source{Lesson 9, videos 5--8; H\&P §3.6.}
\begin{enumerate}
  \item Loads and stores execute in program order among themselves, as in
    Lesson 7.
  \item The load waits until every earlier store has computed its address, so
    that its search never meets an unknown address.
  \item The load executes anyway, and a wrong value is detected and corrected
    later.
\end{enumerate}
The first two options make loads wait; the third does not, but needs a check.
A store's address is unknown until its base register is available, which can
take long when that register is computed by a slow instruction.  In the
following program the address of the store depends on a multiplication,
whereas the address of the load is available almost at once:
\begin{lstlisting}[language={[x86masm]Assembler}, numbers=none]
I1:  MUL R1, R2, R3     ; R1 <- R2 * R3
I2:  ADD R1, R1, R4     ; R1 <- R1 + R4 (an address)
I3:  SW  R5, 0(R1)      ; MEM[R1] <- R5
I4:  ADD R6, R6, R7     ; R6 <- R6 + R7 (an address)
I5:  LW  R8, 0(R6)      ; R8 <- MEM[R6]
I6:  SUB R9, R8, R5     ; R9 <- R8 - R5
\end{lstlisting}
We use the timing conventions of \cref{sec:l08-example}.  MUL takes 4 cycles;
ADD, SUB, loads and stores take 1 cycle, in which a load computes its address
and obtains its value and a store computes its address; the value R5 is
available from the start.  I1 writes its result in cycle~6 and I2 in cycle~8,
so the store I3 executes, and its address becomes known, only in cycle~9.  I4
writes its result in cycle~6, so the load I5 could execute in cycle~7
(\cref{fig:l09-policies}).

\begin{figure}[htbp]
  \centering
  % Out-of-order versus in-order load-store execution for the program of
% section "Out-of-order and in-order load-store execution".  The store I3
% computes its address only in cycle 9; the load I5 could execute in cycle 7.
\begin{tikzpicture}[x=0.52cm,y=0.5cm, font=\scriptsize\sffamily, >={Stealth[length=1.6mm]},
  c/.style={minimum width=0.52cm, minimum height=0.5cm, inner sep=0pt, draw=black},
  iss/.style={c, fill=white},
  ex/.style={c, fill=ln-accent!25},
  sto/.style={c, fill=ln-caution!20},
  wr/.style={c, fill=ln-accent, text=white, font=\scriptsize\bfseries\sffamily},
  hd/.style={font=\scriptsize\bfseries\sffamily, anchor=west}]
  % cycle axes and grids of both panels
  \foreach \ya/\yb/\yc in {0/-0.4/-6.5, -8/-8.4/-11.5} {
    \foreach \c in {1,...,13} { \node[text=ln-muted] at (\c,\ya) {\c}; }
    \node[anchor=east, text=ln-muted] at (0.4,\ya) {\iflangja{サイクル}{cycle}};
    \foreach \c in {1,...,14} { \draw[ln-rule, very thin] (\c-0.5,\yb) -- (\c-0.5,\yc); }
  }
  % ------------------------------------------------ (a) out-of-order
  \node[hd] at (-6.2,0.9) {\iflangja{(a) アウトオブオーダのロード／ストア実行}%
                                     {(a) out-of-order load-store execution}};
  \foreach \y/\lab/\i/\ea/\eb/\w/\st in {
      -1/{I1\enskip\texttt{MUL R1,R2,R3}}/1/2/5/6/ex,
      -2/{I2\enskip\texttt{ADD R1,R1,R4}}/2/7/7/8/ex,
      -3/{I3\enskip\texttt{SW\ \ R5,0(R1)}}/3/9/9/0/sto,
      -4/{I4\enskip\texttt{ADD R6,R6,R7}}/4/5/5/6/ex,
      -5/{I5\enskip\texttt{LW\ \ R8,0(R6)}}/5/7/7/8/ex,
      -6/{I6\enskip\texttt{SUB R9,R8,R5}}/6/9/9/10/ex,
      -9/{I3\enskip\texttt{SW\ \ R5,0(R1)}}/3/9/9/0/sto,
      -10/{I5\enskip\texttt{LW\ \ R8,0(R6)}}/5/10/10/11/ex,
      -11/{I6\enskip\texttt{SUB R9,R8,R5}}/6/12/12/13/ex} {
    \node[anchor=east] at (0.4,\y) {\lab};
    \node[iss] at (\i,\y) {I};
    \pgfmathtruncatemacro{\lnwa}{\i+1}
    \pgfmathtruncatemacro{\lnwb}{\ea-1}
    \ifnum\lnwa<\ea
      \draw[ln-muted, thick] (\lnwa-0.5,\y) -- (\lnwb+0.5,\y);
    \fi
    \foreach \e in {\ea,...,\eb} { \node[\st] at (\e,\y) {E}; }
    \ifnum\w>0 \node[wr] at (\w,\y) {W}; \fi
  }
  % check of the store
  \draw[->, ln-caution, thick] (9.28,-3.25) -- (9.28,-4.3) -- (7.3,-4.3) -- (7.3,-4.72);
  \node[anchor=west, align=left, text=ln-caution, fill=white, inner sep=1pt] at (10.6,-4.5)
    {\iflangja{I3 がアドレスを計算し，\\実行済みの I5 を検査}%
              {I3 computes its address\\and checks I5}};
  % ------------------------------------------------ (b) in-order
  \node[hd] at (-6.2,-7.1) {\iflangja{(b) インオーダのロード／ストア実行}%
                                      {(b) in-order load-store execution}};
  \draw[->, ln-accent, thick] (9.3,-9.25) -- (9.7,-9.75);
  % ------------------------------------------------ legend
  \begin{scope}[shift={(0,-12.6)}]
    \node[iss] at (-5.6,0) {I}; \node[anchor=west] at (-5.1,0) {\iflangja{発行}{issue}};
    \draw[ln-muted, thick] (-2.7,0) -- (-1.7,0);
    \node[anchor=west] at (-1.6,0) {\iflangja{待機}{waiting}};
    \node[ex] at (1.6,0) {E}; \node[anchor=west] at (2.1,0) {\iflangja{実行}{execute}};
    \node[sto] at (5.0,0) {E}; \node[anchor=west] at (5.5,0) {\iflangja{ストア実行}{store executes}};
    \node[wr] at (10.0,0) {W}; \node[anchor=west] at (10.5,0) {\iflangja{結果書き込み}{write result}};
  \end{scope}
\end{tikzpicture}

  \caption{The program of \cref{sec:l09-policies} under two policies.  (a) I5
    executes before the store I3 knows its address; I3 checks I5 in cycle~9.
    (b) I5 waits until I3 has computed its address.  I1, I2 and I4 behave as in
    (a).}
  \label{fig:l09-policies}
\end{figure}

\subsection{Out-of-order load-store execution}

\begin{definition}[Out-of-order load-store execution]
  \label{def:l09-ooo}
  Under \term{out-of-order load-store execution}[ロードストアのアウトオブオーダ実行]
  (option~3), a load executes as soon as its address can be computed, whether
  or not the earlier stores have computed theirs.  Its LSQ search passes over
  stores with unknown addresses as if their addresses differed.
\end{definition}

In the example I5 executes in cycle~7 and I6, which uses its value, in
cycle~9.  In return, every store must check its own address against the loads
that have overtaken it.  When I3 computes its address in cycle~9, it searches
the later LSQ entries for loads with the same address whose C bit is set, that
is, loads that have already obtained a value.
\begin{itemize}
  \item If the address of I5 differs, nothing happens: I5 has obtained the
    correct value, and I6 writes its result in cycle~10.
  \item If the addresses are equal, I5 has obtained a value that I3
    overwrites in program order, so every instruction that used this value,
    directly or indirectly, has computed a wrong result.  The processor cancels
    I5 and all instructions after it and executes them again, starting by
    fetching I5.  None of them has committed, since I3 precedes them, so the
    ROB discards them like the instructions of a mispredicted path
    (\cref{sec:l08-recovery}).  When I5 executes again, its search finds I3,
    and it receives the value by forwarding.
\end{itemize}

\subsection{In-order load-store execution}

\begin{definition}[In-order load-store execution]
  \label{def:l09-in-order}
  Under \term{in-order load-store execution}[ロードストアのインオーダ実行]
  (option~1), loads and stores execute in program order among themselves: a
  load or store waits until every earlier load and store has executed, where a
  store counts as executed once its address is known.  Other instructions
  still execute out of order.
\end{definition}

In the example I5 waits for I3 and executes in cycle~10, and I6 executes in
cycle~12 and writes its result in cycle~13.  No load ever obtains a wrong
value, and nothing is re-executed.  But I5 waits three cycles also when its
address differs from that of I3, and then the wait is pure loss.  Option~2
gives the same timing here, because no load precedes I5; in general it waits
less, since a load does not depend on an earlier load.

\subsection{Comparison}

Out-of-order load-store execution gains cycles whenever a load does not read a
location written by an earlier store whose address is still unknown, and loses
cycles when it does, because cancelling and re-executing the load and the
instructions after it costs more than waiting would have.  Options 1 and 2
delay every load that follows a store with an unknown address, whether the two
conflict or not.  In real programs few loads actually read what such a store
writes, so out-of-order load-store execution performs better, and it is the
policy that high-performance processors use.

\begin{keypoint}
  In-order load-store execution never re-executes a load but makes loads wait
  for stores they usually do not depend on.  Out-of-order load-store execution
  lets loads run at once; correctness then relies on each store checking the
  completed later loads when it computes its address.
\end{keypoint}

\section{Store-to-load forwarding}
\label{sec:l09-forwarding}

Store-to-load forwarding relates each load to the store whose value it must
read, and a processor with out-of-order load-store execution uses this relation
from both ends.\source{Lesson 9, video 9; H\&P ch.~3.}
\begin{description}
  \item[For a load] the value comes from the nearest earlier store with the
    same address, the last store to write the location before the load in
    program order.  An older store to the same address does not qualify: its
    value is overwritten by the nearer one.  If the LSQ holds no earlier store
    with this address, the value comes from memory.\caution{Search from the
    load backwards: a search forward from the commit pointer finds the oldest
    matching store.}
  \item[For a store] the value goes to every later load with the same address,
    up to the next store with that address.  The loads after that store receive
    the value of that store instead.
\end{description}
A load applies the first rule in its search (\cref{sec:l09-lsq}), and a store
applies the second in its check of the completed later loads
(\cref{sec:l09-policies}).  \Cref{fig:l09-forwarding} shows the relation for
seven loads and stores to two addresses.

\begin{figure}[htbp]
  \centering
  % Store-to-load forwarding for a sequence of seven loads and stores: each load
% receives the value of the nearest earlier store with the same address, or
% the value in memory if there is none.  A store's value reaches the later
% loads with its address up to the next store with that address.
\begin{tikzpicture}[x=1cm,y=1cm, font=\footnotesize\sffamily, >={Stealth[length=1.8mm]},
  hd/.style={font=\footnotesize\bfseries\sffamily},
  row/.style={draw, minimum height=0.45cm, inner sep=0pt},
  lab/.style={font=\scriptsize\sffamily, text=ln-muted}]
  % column headings
  \node[hd] at (0.35,0.45) {\iflangja{命令}{Entry}};
  \node[hd] at (1.05,0.45) {L/S};
  \node[hd] at (1.95,0.45) {Addr};
  \node[hd] at (2.9,0.45) {Value};
  \foreach \r/\n/\t/\a/\v/\f in {1/S1/S/0x40/5/ln-accent!12, 2/L2/L/0x48/30/ln-tip!12,
      3/S3/S/0x48/6/ln-tip!12, 4/L4/L/0x40/5/ln-accent!12, 5/S5/S/0x40/9/ln-accent!12,
      6/L6/L/0x48/6/ln-tip!12, 7/L7/L/0x40/9/ln-accent!12} {
    \pgfmathsetmacro{\y}{-(\r-1)*0.45}
    \node[row, fill=ln-box, minimum width=0.7cm, text=ln-muted] at (0.35,\y) {\n};
    \node[row, fill=\f, minimum width=0.7cm] at (1.05,\y) {\t};
    \node[row, fill=\f, minimum width=1.1cm, font=\scriptsize\ttfamily] at (1.95,\y) {\a};
    \node[row, fill=\f, minimum width=0.8cm] at (2.9,\y) {\v};
  }
  % address 0x40 (blue): S1 -> L4, S5 -> L7 (right side)
  \draw[thick, ln-accent] (3.3,0.0) -- (3.75,0.0) -- (3.75,-1.35);
  \draw[->, thick, ln-accent] (3.75,-1.35) -- (3.3,-1.35);
  \draw[thick, ln-accent] (3.3,-1.8) -- (3.75,-1.8) -- (3.75,-2.7);
  \draw[->, thick, ln-accent] (3.75,-2.7) -- (3.3,-2.7);
  \node[font=\scriptsize\sffamily, text=ln-accent, anchor=west, align=left] at (3.95,-0.675)
    {\iflangja{S1 の値：\\S5 の手前まで}{value of S1:\\up to S5}};
  \node[font=\scriptsize\sffamily, text=ln-accent, anchor=west, align=left] at (3.95,-2.25)
    {\iflangja{S5 の値}{value of S5}};
  % address 0x48 (green): memory -> L2, S3 -> L6 (left side)
  \node[draw, fill=white, font=\scriptsize\sffamily, inner sep=2pt] (mem) at (-1.35,-0.45)
    {\iflangja{メモリ}{memory}};
  \draw[->, thick, ln-tip!80!black] (mem.east) -- (0.0,-0.45);
  \draw[thick, ln-tip!80!black] (0.0,-0.9) -- (-0.45,-0.9) -- (-0.45,-2.25);
  \draw[->, thick, ln-tip!80!black] (-0.45,-2.25) -- (0.0,-2.25);
  \node[font=\scriptsize\sffamily, text=ln-tip!80!black, anchor=east, align=right] at (-0.6,-1.575)
    {\iflangja{S3 の値}{value of S3}};
  % program order
  \draw[->, ln-muted] (5.75,0.2) -- (5.75,-2.9);
  \node[lab, rotate=-90] at (6.0,-1.35) {\iflangja{プログラム順}{program order}};
\end{tikzpicture}

  \caption{Store-to-load forwarding for seven loads and stores to
    \texttt{0x40} and \texttt{0x48}, with the values of a sequential
    execution.  L4 receives the value of S1, not of S5, which follows it; L7
    receives the value of S5, which overwrites that of S1.  L2 precedes S3 and
    reads memory.}
  \label{fig:l09-forwarding}
\end{figure}

\section{A load-store queue example}
\label{sec:l09-example}

We now trace the loads and stores of \cref{fig:l09-forwarding}
(\cref{ex:l09-lsq}) under out-of-order load-store execution.\source{Lesson 9,
video 10; H\&P §3.6.}

\begin{example}
  \label{ex:l09-lsq}
  The loads and stores of a program are, in program order,
  \begin{lstlisting}[language={[x86masm]Assembler}, numbers=none]
S1:  SW  R1, 0(R2)     ; MEM[0x40] <- R1
L2:  LW  R3, 8(R2)     ; R3 <- MEM[0x48]
S3:  SW  R4, 8(R10)    ; MEM[R10+8] <- R4
L4:  LW  R5, 0(R2)     ; R5 <- MEM[0x40]
S5:  SW  R6, 0(R7)     ; MEM[R7] <- R6
L6:  LW  R8, 8(R2)     ; R8 <- MEM[0x48]
L7:  LW  R9, 0(R2)     ; R9 <- MEM[0x40]
\end{lstlisting}
  with $\text{R2} = \texttt{0x40}$, $\text{R1}=5$, $\text{R4}=6$ and
  $\text{R6}=9$ available from the start.  $\text{R10} = \texttt{0x40}$ and
  $\text{R7}=\texttt{0x40}$ are computed by slow instructions that precede S3
  and S5, so these two stores learn their addresses late, in the order of
  \cref{tab:l09-trace}.  Memory initially holds 20 at \texttt{0x40} and 30 at
  \texttt{0x48}.
\end{example}

\begin{table}[htbp]
  \centering\small
  \caption{Execution of \cref{ex:l09-lsq}.  In each step one entry computes
    its address: a load searches the earlier stores; a store checks the later
    loads whose C bit is set.}
  \label{tab:l09-trace}
  \begin{tabular}{@{}clcll@{}}
    \toprule
    Step & Entry & Addr & Search or check & Outcome \\
    \midrule
    1 & S1 & \texttt{0x40} & no later load completed & value 5 recorded \\
    2 & L2 & \texttt{0x48} & S1: other address & memory: 30 \\
    3 & L4 & \texttt{0x40} & S3: unknown; S1: match & from S1: 5 \\
    4 & S3 & \texttt{0x48} & L4: other address & no conflict \\
    5 & L6 & \texttt{0x48} & S5: unknown; S3: match & from S3: 6 \\
    6 & L7 & \texttt{0x40} & S5: unknown; S3: other; S1: match & from S1: 5 \\
    7 & S5 & \texttt{0x40} & L6: other address; L7: match & L7 re-executed \\
    8 & L7 & \texttt{0x40} & S5: match & from S5: 9 \\
    \bottomrule
  \end{tabular}
\end{table}

L4 and L6 pass over stores whose addresses are still unknown, and the later
checks of S3 and S5 find no conflict with them.  L7 also passes over S5, but in
step~7 S5 computes \texttt{0x40} and finds L7 completed with this address, with
no store to \texttt{0x40} between them.  The value 5 that L7 obtained from S1
is overwritten by S5 in program order, so L7 and every instruction after it are
cancelled and executed again, and in step~8 L7 receives 9 from S5.

After the entries commit in program order, memory holds 9 at \texttt{0x40}
and 6 at \texttt{0x48}, and the loads have delivered $\text{R3}=30$,
$\text{R5}=5$, $\text{R8}=6$ and $\text{R9}=9$, the results of a sequential
execution.  With in-order load-store execution L4 would have waited for S3,
and L6 and L7 for S5, and nothing would have been executed twice.

\section{The LSQ, the ROB and the reservation stations}
\label{sec:l09-rob-rs}

Loads and stores pass through the same steps as other instructions, and the LSQ
adds work to each step (\cref{tab:l09-steps}).\source{Lesson 9, videos 11--14;
H\&P §3.6.}  A load or store is issued only when a ROB entry and an LSQ entry
are both available; the LSQ entry takes the place of the reservation station
that other instructions need with their ROB entry.  A load broadcasts its value on the CDB under the name of its ROB entry, exactly
like an arithmetic result.  A store produces no register value, so nothing is
broadcast, and its ROB entry is done as soon as the store has recorded both its
address and its value.

\begin{table}[htbp]
  \centering\small
  \caption{The steps of a load and a store with an LSQ.}
  \label{tab:l09-steps}
  \begin{tabular}{@{}>{\raggedright\arraybackslash}p{0.15\linewidth}>{\raggedright\arraybackslash}p{0.38\linewidth}>{\raggedright\arraybackslash}p{0.38\linewidth}@{}}
    \toprule
    Step & Load & Store \\
    \midrule
    Issue & needs a free ROB entry and LSQ entry (no RS); the base register is
      looked up in the RAT, and the RAT entry of the destination register is
      set to the load's ROB entry & needs a free ROB entry and LSQ entry (no
      RS); the base and value registers are looked up in the RAT; no RAT entry
      changes \\
    \addlinespace
    Execute & computes the address and records it in the LSQ; searches the
      earlier stores; copies a forwarded value or reads memory; sets C
      & records the address when the base register is available and checks
      the later loads whose C is set; records the value when it is available;
      sets C once both are recorded \\
    \addlinespace
    Write result & broadcasts the value on the CDB; the ROB entry receives it
      & nothing to broadcast \\
    \addlinespace
    Commit & the value is written to the register file, and the RAT entry is
      reset if it still names this ROB entry; the ROB and LSQ entries are
      freed & the value is written to memory; the ROB and LSQ entries are
      freed \\
    \bottomrule
  \end{tabular}
\end{table}

Both queues are in program order, so the load or store at the commit pointer of
the ROB is also the oldest entry of the LSQ, and the two entries are freed
together.\index{commit}  A store leaves the LSQ exactly when its value is
written to memory, so a later load always finds that value in one of the two
places it looks: in the store's LSQ entry before the commit, in memory after
it.

\begin{keypoint}[Lesson summary]
  Loads and stores depend on one another through memory addresses that are
  known only after execution, and memory is not renamed.  Stores therefore
  write memory at commit, in program order.  The load-store queue holds the
  uncommitted loads and stores in program order, with their addresses, values
  and completion bits.  A load takes its value from the nearest earlier store
  with the same address by store-to-load forwarding, or from memory if there
  is none.  Earlier stores with unknown addresses
  are handled either by in-order load-store execution, where the load waits for
  them, or by out-of-order load-store execution, where the load proceeds and a
  store that later computes a matching address has the load and all
  instructions after it re-executed; the latter is faster because such
  conflicts are rare.  A load or store needs a ROB entry and an LSQ entry at
  issue, where other instructions need a ROB entry and an RS.
\end{keypoint}

\end{document}
